\documentclass[10pt,xcolor=dvipsnames]{beamer}
\usetheme{Madrid}
\usecolortheme[RGB={0,82,155}]{structure}
\setbeamertemplate{navigation symbols}{}

%% ── Idiomas y fuentes ──
\usepackage[utf8]{inputenc}
\usepackage[spanish]{babel}
\usepackage{amsmath,amssymb}
\usepackage{hyperref}

%% ── Portada ──
\title[Relaciones y Funciones]{Relaciones y Funciones}
\subtitle{Matemática Discreta --- MM420 \\ \small Capítulo 5}
\author[L. Berríos]{Luis Berríos}
\institute[UNAH]{Universidad Nacional Autónoma de Honduras \\ Facultad de Ciencias \\ Departamento de Matemáticas}
\date{II Período, 2026}

%% ── Índice al inicio ──
\AtBeginSection[]{%
  \begin{frame}<beamer>\frametitle{Contenido}
    \tableofcontents[currentsection]
  \end{frame}
}

%% ======================================================================
\begin{document}
%% ======================================================================

\begin{frame}
  \titlepage
\end{frame}

\begin{frame}{Contenido}
  \tableofcontents
\end{frame}





%% ======================================================================
\section{Productos cartesianos y relaciones}
%% ======================================================================

\begin{frame}{Producto cartesiano}
  \begin{block}{Definición}
    Sean $A$ y $B$ conjuntos. El \textbf{producto cartesiano} $A \times B$ es el
    conjunto de todos los \alert{pares ordenados} $(a,b)$ con $a \in A$ y $b \in B$:
    \[
      A \times B = \{(a,b) \mid a \in A \text{ y } b \in B\}.
    \]
  \end{block}

  \begin{itemize}
    \item Si $A$ y $B$ son finitos, $|A \times B| = |A| \cdot |B|$.
    \item El par es \textbf{ordenado}: $(a,b) \neq (b,a)$ en general.
  \end{itemize}
\end{frame}

\begin{frame}{Relación}
  \begin{block}{Definición}
    Una \textbf{relación} de $A$ en $B$ es un subconjunto de $A \times B$.
    Una \textbf{relación en} $A$ es un subconjunto de $A \times A$.
  \end{block}

  \begin{block}{Notación}
    Si $R \subseteq A \times B$, escribimos $a \mathbin{R} b$ (ó $a\,R\,b$) en
    lugar de $(a,b) \in R$.
  \end{block}

  \begin{itemize}
    \item Ejemplo: en $\mathbb{Z}$, ``$a \leq b$'' es una relación.
    \item Una relación queda determinada por su \textbf{grafo}: $\{(a,b) \in R\}$.
  \end{itemize}
\end{frame}





%% ======================================================================
\section{Funciones: en general e inyectivas}
%% ======================================================================

\begin{frame}{Función}
  \begin{block}{Definición}
    Una \textbf{función} $f$ de $A$ en $B$, escrita $f : A \to B$, es una
    relación de $A$ en $B$ tal que \alert{cada elemento de $A$ aparece como
    primera componente de exactamente un par}.
  \end{block}

  \begin{itemize}
    \item $A$ es el \textbf{dominio} de $f$; $B$ es el \textbf{codominio}.
    \item Si $f(a) = b$, $b$ es la \textbf{imagen} de $a$, y $a$ es una
          \textbf{preimagen} de $b$.
    \item El \textbf{rango} (o imagen) de $f$ es
          $f(A) = \{f(a) \mid a \in A\} \subseteq B$.
  \end{itemize}
\end{frame}

\begin{frame}{Función inyectiva}
  \begin{block}{Definición}
    Una función $f : A \to B$ es \textbf{inyectiva} (o \emph{uno a uno}) si
    \alert{elementos distintos de $A$ se llevan a elementos distintos de $B$}:
    \[
      \forall a_1, a_2 \in A,\; f(a_1) = f(a_2) \implies a_1 = a_2.
    \]
  \end{block}

  \begin{itemize}
    \item Equivalentemente: $a_1 \neq a_2 \implies f(a_1) \neq f(a_2)$.
    \item Si $A$ es finito y $f$ es inyectiva, entonces $|A| \leq |B|$.
  \end{itemize}
\end{frame}





%% ======================================================================
\section{Funciones sobreyectivas: Números de Stirling}
%% ======================================================================

\begin{frame}{Función sobreyectiva}
  \begin{block}{Definición}
    Una función $f : A \to B$ es \textbf{sobreyectiva} (o \emph{subyectiva}) si
    \alert{todo elemento de $B$ es imagen de algún elemento de $A$}:
    \[
      \forall b \in B,\; \exists a \in A \text{ tal que } f(a) = b.
    \]
    Es decir, $f(A) = B$.
  \end{block}

  \begin{itemize}
    \item Si $A$ y $B$ son finitos y $f : A \to B$ es sobreyectiva,
          entonces $|A| \geq |B|$.
    \item \textbf{Biyectiva}: inyectiva \emph{y} sobreyectiva.
  \end{itemize}
\end{frame}

\begin{frame}{Conteo de funciones sobreyectivas}
  \begin{block}{Número de funciones sobreyectivas}
    Para $A, B$ finitos con $|A| = n$ y $|B| = k$ ($n \geq k$), el número de
    funciones sobreyectivas $A \to B$ es
    \[
      k! \cdot S(n,k),
    \]
    donde $S(n,k)$ es el \textbf{número de Stirling de segundo tipo}.
  \end{block}
\end{frame}

\begin{frame}{Números de Stirling del segundo tipo}
  \begin{block}{Definición}
    Para enteros $n, k \geq 0$, el \textbf{número de Stirling de segundo tipo}
    $S(n,k)$ es el \alert{número de formas de particionar un conjunto de
    $n$ elementos en $k$ subconjuntos no vacíos}.
  \end{block}

  \begin{itemize}
    \item Convención: $S(n,0) = 0$ si $n > 0$, y $S(0,0) = 1$.
    \item \textbf{Recurrencia}:
      \[
        S(n,k) = k \cdot S(n-1,k) + S(n-1,k-1), \quad 1 \leq k \leq n.
      \]
    \item Primeros valores: $S(n,1) = S(n,n) = 1$;
          $S(3,2) = 3$, $S(4,2) = 7$, $S(4,3) = 6$.
  \end{itemize}
\end{frame}





%% ======================================================================
\section{Funciones especiales}
%% ======================================================================

\begin{frame}{Funciones especiales}
  \begin{block}{Función constante}
    $f : A \to B$ es \textbf{constante} si existe $b_0 \in B$ tal que
    $f(a) = b_0$ para todo $a \in A$.
  \end{block}

  \begin{block}{Función identidad}
    La \textbf{identidad} en $A$ es $\mathrm{id}_A : A \to A$ definida por
    $\mathrm{id}_A(a) = a$ para todo $a \in A$. Es biyectiva.
  \end{block}

  \begin{block}{Función piso y techo}
    Para $x \in \mathbb{R}$:
    \[
      \lfloor x \rfloor = \text{mayor entero } \leq x, \qquad
      \lceil x \rceil = \text{menor entero } \geq x.
    \]
  \end{block}
\end{frame}

\begin{frame}{Funciones especiales: lógica y factorial}
  \begin{block}{Función característica}
    Para un subconjunto $X \subseteq U$,
    \[
      \chi_X : U \to \{0,1\}, \quad
      \chi_X(u) = \begin{cases} 1 & \text{si } u \in X, \\ 0 & \text{si } u \notin X. \end{cases}
    \]
  \end{block}

  \begin{block}{Factorial y combinaciones}
    Para $n \geq 0$,
    $n! = n \cdot (n-1) \cdots 1$ (con $0! = 1$).
    Para $0 \leq r \leq n$,
    \[
      \binom{n}{r} = \frac{n!}{r!\,(n-r)!}
    \]
    cuenta los subconjuntos de tamaño $r$ de un conjunto de $n$ elementos.
  \end{block}
\end{frame}





%% ======================================================================
\section{El principio del palomar}
%% ======================================================================

\begin{frame}{Principio del palomar}
  \begin{block}{Enunciado}
    Si se distribuyen $n$ objetos en $k$ cajas y $n > k$, entonces
    \alert{al menos una caja contiene más de un objeto}.
  \end{block}

  \begin{itemize}
    \item \textbf{Versión general}: si $n$ objetos van a $k$ cajas, alguna caja
          contiene al menos $\lceil n/k \rceil$ objetos.
    \item Es \emph{no constructivo}: garantiza la existencia, no la exhibe.
  \end{itemize}

  \begin{exampleblock}{Aplicaciones típicas}
    Demostrar que entre ciertos conjuntos existen dos elementos con la misma
    propiedad (misma fecha de cumpleaños, mismo residuo módulo $k$, etc.).
  \end{exampleblock}
\end{frame}





%% ======================================================================
\section{Composición de funciones e inversas}
%% ======================================================================

\begin{frame}{Composición de funciones}
  \begin{block}{Definición}
    Dadas $g : A \to B$ y $f : B \to C$, la \textbf{composición} $f \circ g :
    A \to C$ se define por
    \[
      (f \circ g)(a) = f(g(a)), \quad a \in A.
    \]
  \end{block}

  \begin{itemize}
    \item La composición \emph{no} es conmutativa en general.
    \item \textbf{Identidad}: $f \circ \mathrm{id}_A = f = \mathrm{id}_B \circ f$.
    \item \textbf{Asociatividad}: $(f \circ g) \circ h = f \circ (g \circ h)$.
    \item Composición de inyectivas es inyectiva; de sobreyectivas es sobreyectiva.
  \end{itemize}
\end{frame}

\begin{frame}{Función inversa}
  \begin{block}{Definición}
    Si $f : A \to B$ es \textbf{biyectiva}, su \textbf{inversa} $f^{-1} : B \to A$
    es la función definida por
    \[
      f^{-1}(b) = a \iff f(a) = b.
    \]
  \end{block}

  \begin{itemize}
    \item $f \circ f^{-1} = \mathrm{id}_B$ y $f^{-1} \circ f = \mathrm{id}_A$.
    \item $f$ es biyectiva $\iff$ tiene inversa.
    \item Si $f$ y $g$ son biyectivas, $(g \circ f)^{-1} = f^{-1} \circ g^{-1}$.
  \end{itemize}
\end{frame}





%% ======================================================================
\section{Complejidad computacional}
%% ======================================================================

\begin{frame}{Complejidad computacional}
  \begin{block}{Idea}
    La \textbf{complejidad computacional} estudia los recursos (tiempo,
    memoria) que requiere un algoritmo en función del tamaño de la entrada.
  \end{block}

  \begin{itemize}
    \item Un \textbf{problema computacional} es una pregunta genérica a
          resolver; un \textbf{algoritmo} es un procedimiento paso a paso.
    \item \textbf{Instancia}: un caso particular del problema (con datos
          específicos).
  \end{itemize}

  \begin{block}{Notación asintótica (repaso)}
    $f(n) = O(g(n))$ si existen $C, n_0 > 0$ tales que
    $|f(n)| \leq C \cdot g(n)$ para todo $n \geq n_0$.
  \end{block}
\end{frame}

\begin{frame}{Clases de complejidad usuales}
  \begin{block}{Jerarquía de crecimiento}
    De menor a mayor:
    \[
      1 \;<\; \log n \;<\; n \;<\; n \log n \;<\; n^2 \;<\; n^3 \;<\; 2^n.
    \]
  \end{block}

  \begin{itemize}
    \item \textbf{Tiempo polinomial}: $O(n^k)$ para alguna constante $k$
          (clase \textbf{P}).
    \item \textbf{Tiempo exponencial}: $O(c^n)$ con $c > 1$.
  \end{itemize}
\end{frame}





%% ======================================================================
\section{Análisis de algoritmos}
%% ======================================================================

\begin{frame}{Análisis de algoritmos}
  \begin{block}{Definición}
    El \textbf{análisis de algoritmos} determina, asintóticamente, el número de
    operaciones elementales que un algoritmo realiza en el peor caso (o caso
    promedio) en función del tamaño $n$ de la entrada.
  \end{block}

  \begin{itemize}
    \item \textbf{Peor caso}: cota superior para \emph{todas} las entradas de
          tamaño $n$.
    \item \textbf{Caso promedio}: esperado sobre una distribución de entradas.
    \item \textbf{Mejor caso}: usualmente poco útil, pero informativo.
  \end{itemize}

  \begin{exampleblock}{Ejemplo}
    Búsqueda lineal en un arreglo de tamaño $n$:
    peor caso $T(n) = O(n)$.
  \end{exampleblock}
\end{frame}





%% ======================================================================
\section{Resumen y repaso histórico}
%% ======================================================================

\begin{frame}{Resumen del capítulo}
  \begin{itemize}
    \item \textbf{Producto cartesiano}: pares ordenados; $|A \times B| = |A||B|$.
    \item \textbf{Relación}: subconjunto de un producto cartesiano.
    \item \textbf{Función}: relación con primera componente única.
    \item \textbf{Inyectiva / sobreyectiva / biyectiva}: tipos de funciones.
    \item \textbf{Stirling de segundo tipo} $S(n,k)$: particiones en $k$ bloques;
          funciones sobreyectivas $= k!\,S(n,k)$.
    \item \textbf{Constante, identidad, piso, techo, característica, factorial}.
    \item \textbf{Principio del palomar}: $n$ objetos en $k$ cajas $\Rightarrow$
          una caja tiene $\geq \lceil n/k \rceil$ objetos.
    \item \textbf{Composición} $f \circ g$ y \textbf{inversa} $f^{-1}$.
    \item \textbf{Complejidad} y \textbf{análisis de algoritmos}: $O(\cdot)$,
          $P$, polinomial vs.\ exponencial.
  \end{itemize}
\end{frame}

\begin{frame}{Repaso histórico}
  \begin{itemize}
    \item \textbf{Georg Cantor} (s.\ XIX): fundador de la teoría de conjuntos.
    \item \textbf{James Stirling} (s.\ XVIII): números que llevan su nombre,
          con aplicaciones en combinatoria.
    \item \textbf{Dirichlet} (s.\ XIX): formulaciones del principio del palomar
          (\emph{cajón de Dirichlet}).
    \item \textbf{Alonzo Church} y \textbf{Alan Turing} (s.\ XX): bases
          teóricas de la computación y de la noción de algoritmo.
    \item \textbf{Don Knuth} (s.\ XX): ``El arte de programar'', formaliza el
          análisis asintótico de algoritmos.
  \end{itemize}
\end{frame}

\end{document}
